% songs/eine-insel-mit-zwei-bergen.tex
\begin{song}{Eine Insel Mit Zwei Bergen}{Augsburger Puppenkiste}{Capo -}

\cl{  G  D  G}
\slbl{1.}Ei-ne Insel mit zwei Bergen
\cl{       D}
und im tiefen, weiten Meer
\cl{         D7}
Mit viel Tunnels und Geleisen
\cl{        G}
und dem Eisenbahnverkehr

Nun, wie mag die Insel heißen,
\cl{        D}
ringshe-rum ist schöner Strand
\cl{      D7}
Jeder sollte einmal reisen,
\cl{         G}
In das schöne Lummerland!
\cl{   Am | G | D | G}
Lalal...

\cl{  G  D  G}
\slbl{2.}Ei-ne Insel mit zwei Bergen
\cl{        D}
Und dem Fotoatelier
\cl{       D7}
In dem letzten macht man Bilder,
\cl{        G}
Auf den ersten "dulijö"
\cl{G}
Diese Breiten, diese Tiefen,
\cl{      D}
Diese Höhen sind bekannt
\cl{        D7}
Und man spricht von den Motiven
\cl{        G}
Auf dem schönen Lummerland!
\cl{   Am | G | D | G}
Lalal...

\cl{  G  D  G}
\slbl{3.}Eine Insel mit zwei Bergen
\cl{        D}
Und dem Fernsprechtelefon 
\cl{          D7}
wählt man nur die richtige Nummer,
\cl{                G}
klappt auch die Verbindung schon.
\cl{G}
"Hallo, hier ist falsch verbunden!"
\cl{        D}
"Wollen sie sich jetzt beschwer´n?"
\cl{       D7}
"Nein, warum? Das kann passier´n!"
\cl{       G}
"Also dann auf Wiederhör´n!"
\cl{   Am | G | D | G}
Lalal...

\cl{  G  D  G}
\slbl{4.}Eine Insel mit zwei Bergen
\cl{        D}
Und dem Laden von Frau Waas:
\cl{      D7}
Hustenbonbons, Alleskleber
\cl{      G}
Regen-schirme, Leberkas,
\cl{G}
Körbe, Hüte, Lampen, Würste,
\cl{       D}
Blumen-kohl und Fensterglas,
\cl{      D7}
Leder-hosen, Kukucksuhren
\cl{         G}
Und noch dies und dann noch das!
\cl{   Am | G | D | G}
Lalal...
\anmerkung{Jim Knopf und Lukas der Lokomotivführer gilt als Gegengeschichte zu nationalsozialistischen Bilderwelten und zeitgenössischer Fehlinterpretationen des Darwinismus.}
\griffe{D7}
\end{song}
